% ============ 第3章 §3.3 Lebesgue 微分定理 (原书 pp.71-74) ============

\section{Lebesgue 微分定理}\label{sec:3.3}
\index{0symbolBr@$B(r)$}

\index{0symbolLloc@$L_{\mathrm{loc}}^{1}$}

\index{L@Lebesgue 点}

\index{L@Lebesgue 微分定理}










前面说过覆盖引理对函数可微性的分析是有用的, 这里说明这一侧面. 这里考察当 $r \downarrow 0$ 时,
\[
\frac{1}{|Q(x,r)|}\int_{Q(x,r)}f(y)\dd y
\]
会呈现怎样的行为. 这里, $Q(x,r)$ 是以 $x$ 为中心、体积为 $(2r)^{n}$ 的立方体.

实际上, 因为想讨论的是与各点收敛相近的概念——几乎处处收敛, 虽然有些离题, 还是重新收录 Lebesgue 收敛定理.

\begin{theorem}[Lebesgue 收敛定理 I]\label{thm:3.6}
设 $\mathbb{R}^{n}$ 中的函数列 $\{f_{j}\}_{j=1}^{\infty}$ 几乎处处收敛于 $f$. 进而, 假设存在 $g \in L^{1}(\mathbb{R}^{n})$, 使得对一切自然数 $j$, 几乎处处有 $|f_{j}| \leq g$. 此时,
\[
\lim_{j \to \infty}\int_{\mathbb{R}^{n}}f_{j}(x)\dd x = \int_{\mathbb{R}^{n}}f(x)\dd x
\]
成立.
\end{theorem}

这一结论用 $L^{1}$ 范数也可以改写如下. 实际上, 这是因为下面的定理与最初的定理等价.

\begin{theorem}[Lebesgue 收敛定理 II]\label{thm:3.7}
设 $\mathbb{R}^{n}$ 中的函数列 $\{f_{j}\}_{j=1}^{\infty}$ 几乎处处收敛于 $f$. 进而, 假设存在 $g \in L^{1}(\mathbb{R}^{n})$, 使得对一切自然数 $j$, 几乎处处有 $|f_{j}| \leq g$. 此时, $f_{j}$ 在 $L^{1}$ 的拓扑下收敛于 $f$. 即
\[
\lim_{j \to \infty}\|f_{j} - f\|_{L^{1}} = 0
\]
成立.
\end{theorem}

进而, 这个定理的假设可以改写如下.

\begin{theorem}[Lebesgue 收敛定理 III]\label{thm:3.8}
设 $\mathbb{R}^{n}$ 中的函数列 $\{f_{j}\}_{j=1}^{\infty}$ 几乎处处收敛于 $f$. 此时, 若 $\sup_{j \in \mathbb{N}}|f_{j}| \in L^{1}(\mathbb{R}^{n})$, 则 $f_{j}$ 在 $L^{1}$ 的拓扑下收敛于 $f$. 即
\[
\lim_{j \to \infty}\|f_{j} - f\|_{L^{1}} = 0
\]
成立.
\end{theorem}

这样改写后, 就能理解极大算子承担着自然的角色. 也就是说, 作为 Lebesgue 收敛定理 I, II 中的函数 $g$, 它取的是自然的(极小的)对象.

$L_{\mathrm{loc}}^{1}(\mathbb{R}^{n})$ 是指乘以 $\chi_{K}$ 后成为可积函数的函数全体所组成的集合, 其中 $K$ 为任意紧集. 这样的函数称为局部可积函数.

\begin{theorem}[Lebesgue 微分定理]\label{thm:3.9}
设 $f \in L_{\mathrm{loc}}^{1}(\mathbb{R}^{n})$. 此时, 对几乎所有 $x \in \mathbb{R}^{n}$,
\begin{equation}\tag{3.14}
\lim_{r \downarrow 0}\frac{1}{|Q(x,r)|}\int_{Q(x,r)}|f(y) - f(x)|\dd y = 0
\label{eq:3.14}
\end{equation}
成立. 特别地, 对这样的点, $|f(x)| \leq Mf(x)$ 成立.
\end{theorem}

称使式~\eqref{eq:3.14} 成立的点为 Lebesgue 点.

\begin{proof}
后半部分的主张由极大算子的定义是显然的, 因此证明前半部分. 为此, 只需证明使式~\eqref{eq:3.14} 不成立的点的全体所组成的集合的测度为 $0$. 将使式~\eqref{eq:3.14} 不成立的点的全体按水平来考虑. 即, 对一切自然数 $k \in \mathbb{N}$, 证明
\[
E_{k} \equiv \left\{x \in \mathbb{R}^{n}: \limsup_{r \downarrow 0}\frac{1}{|Q(x,r)|}\int_{Q(x,r)}|f(y) - f(x)|\dd y > \frac{1}{k}\right\}
\]
的测度为 $0$ 就足够了. 取 $r > 0$, 令 $B(r) = \{|x| < r\}$. 由于式~\eqref{eq:3.14} 是否成立是局部的问题, 用 $f \cdot \chi_{B(r)}$ 代替 $f$ 来考虑, 可以假设 $f \in L^{1}(\mathbb{R}^{n})$. 进而, 注意对 $g \in L^{1}(\mathbb{R}^{n}) \cap \mathrm{BC}(\mathbb{R}^{n})$, 式~\eqref{eq:3.14} 是显然的\footnote{BC 是有界连续 (bounded continuous) 之意.}. 任意取 $\varepsilon > 0$. 通过证明 $|E_{k}| \leq 2(3^{n}+1)k\varepsilon$, 完成定理的证明.

取 $g \in L^{1}(\mathbb{R}^{n}) \cap \mathrm{BC}(\mathbb{R}^{n})$, 使得 $\|f - g\|_{L^{1}} \leq \varepsilon$. 这由稠密性是可能的. 略记 $H(x, y) = |f(y) - g(y) - f(x) + g(x)|$. 作为集合的包含关系
\begin{equation}
\begin{aligned}
E_{k} &= \left\{x \in \mathbb{R}^{n}: \limsup_{r \downarrow 0}\frac{1}{|Q(x,r)|}\int_{Q(x,r)}|H(x,y)|\dd y > \frac{1}{k}\right\} \\
&\subset \left\{x \in \mathbb{R}^{n}: M[f-g](x) + |f(x) - g(x)| > \frac{1}{k}\right\} \\
&\subset \left\{x \in \mathbb{R}^{n}: M[f-g](x) > \frac{1}{2k}\right\} \cup \left\{x \in \mathbb{R}^{n}: |f(x) - g(x)| > \frac{1}{2k}\right\}
\end{aligned}\tag{3.15}
\label{eq:3.15}
\end{equation}
成立, 因此由测度的次可加性、Chebyshev 不等式与弱 (1,1) 极大不等式,
\[
|E_{k}| \leq 2(3^{n}+1)k\|f - g\|_{L^{1}} \leq 2(3^{n}+1)k\varepsilon
\]
成立. 由于 $\varepsilon > 0$ 是任意的, 式~\eqref{eq:3.14} 随之成立.
\end{proof}

$C_{\mathrm{c}}^{\infty}(\mathbb{R}^{n})$ 是指 $C^{\infty}(\mathbb{R}^{n})$ 级函数中, 在有界集合之外取 $0$ 的函数全体所组成的线性空间.

\begin{note}\label{note:3.2}
与这里处理的 Lebesgue 微分定理相同, 对 $\varphi \in C_{\mathrm{c}}^{\infty}(\mathbb{R}^{n})$ 与 $f \in L_{\mathrm{loc}}^{1}(\mathbb{R}^{n})$, 将关于 $x \in \mathbb{R}^{n}$ 的函数定义为 $\varphi_{t}(x) = t^{-n}\varphi(t^{-1}x)$ 时, a.e. 有
\[
\varphi_{t} * f(x) \to \int_{\mathbb{R}^{n}}\varphi(y)\dd y \times f(x) \quad (t \downarrow 0)
\]
成立. 这也常常被称为 Lebesgue 微分定理. 为证明它, 注意下面两点. 一点是当 $f \in C_{\mathrm{c}}^{\infty}(\mathbb{R}^{n})$ 时结论对一切 $x \in \mathbb{R}^{n}$ 正确, 另一点是由于 $\varphi \in C_{\mathrm{c}}^{\infty}(\mathbb{R}^{n})$ 是具有紧支集的有界函数,
\[
\sup_{t > 0}|\varphi_{t} * f(x)| \lesssim Mf(x) \quad (x \in \mathbb{R}^{n})
\]
成立. 注意这两点, 就能用与定理~\ref{thm:3.9} 相同的方法证明.
\end{note}

\begin{exercise}\label{exer:3.7}
对下面的开区间 $I$ 与函数 $\varphi$, 令 $f(t) = \varphi(t)\chi_{I}(t)$. 另外, 取 $a \in \mathbb{R}$, 令 $F_{a}(t) = \int_{a}^{t}f(s)\dd s$. 在下面的 5 种情形下, 求出 $F_{a}$, 并判定 $F_{a}$ 在 $t = a$ 处是否可微.
\[
\begin{array}{l}
\bullet\ a < \inf I \\
\bullet\ a = \inf I \\
\bullet\ a \in I \\
\bullet\ a = \sup I \\
\bullet\ a > \sup I
\end{array}
\]
(1) $I = (5, 8)$, $\varphi(t) = t - 5$.

(2) $I = (6, 9)$, $\varphi(t) = t - 9$.
\end{exercise}

\begin{exercise}\label{exer:3.8}
回答下面的问题.

(1) 设 $f \in L_{\mathrm{loc}}^{1}(\mathbb{R})$. 此时, 证明对几乎所有 $x \in \mathbb{R}$,
\begin{equation}\tag{3.16}
\lim_{r \downarrow 0}\frac{1}{r}\int_{x}^{x+r}f(y)\dd y = f(x)
\label{eq:3.16}
\end{equation}
成立.

(2) 设 $f \in L_{\mathrm{loc}}^{1}(\mathbb{R}^{n})$. 此时, 证明对几乎所有 $x \in \mathbb{R}^{n}$,
\begin{equation}\tag{3.17}
\lim_{r \downarrow 0}\frac{1}{|B(x,r)|}\int_{B(x,r)}f(y)\dd y = f(x)
\label{eq:3.17}
\end{equation}
成立.
\end{exercise}

\begin{exercise}\label{exer:3.9}
设 $f \in L_{\mathrm{loc}}^{1}(\mathbb{R})$, 证明以下命题.

(1) 当 $x$ 是 $f$ 的 Lebesgue 点时,
\[
\lim_{r \downarrow 0}\frac{1}{|B(x,r)|}\int_{B(x,r)}|f(y) - f(x)|\dd y = 0
\]
成立.

(2) 反之, 使
\[
\lim_{r \downarrow 0}\frac{1}{|B(x,r)|}\int_{B(x,r)}|f(y) - f(x)|\dd y = 0
\]
成立的 $x \in \mathbb{R}$ 是 $f$ 的 Lebesgue 点.
\end{exercise}

\begin{exercise}\label{exer:3.10}
证明不存在 Lebesgue 可测集 $E$, 使得对一切开区间 $I$ 有 $|E \cap I| = \frac{1}{2}|I|$.
\end{exercise}
